\section{模型假设}

为建立合理的数学模型，根据题目条件和工程实际，提出以下基本假设：

\textbf{假设一：}电氢氨装置整体同步运行。碱性电解槽、质子交换膜电解槽和合成氨装置三者在运行调节时保持同步，即同时开停机、同比例调节功率。该假设基于题目"电氢氨装置只有全额开机和停机"的表述，反映了实际工程中氢气产出与合成氨消耗需要匹配的物理约束。

\textbf{假设二：}制氢实际耗电按效率折算。附件5标注"50 kWh/kg（未考虑效率）"，碱性电解槽效率为70\%，PEMEL效率为80\%。实际耗电=理论耗电/效率，碱性实际耗电为50/0.7=71.43 kWh/kgH$_2$，对应10 MW额定功率（140$\times$71.43=10000 kW=10 MW），与题目参数一致。

\textbf{假设三：}产能扩容时设备功率线性同步提升。题目明确"配套电氢氨装置的额定功率将随产能规模呈线性同步提升"，72吨/日产能为36吨/日的两倍，因此所有设备功率翻倍。该假设保证了设备参数随产能变化的一致性。

\textbf{假设四：}各时段内设备运行视作稳态且功率恒定。题目规定"运行分析时段为1小时"，考虑到储能系统、碱性电解槽与PEMEL的电化学响应时间以及合成氨装置的热力学时间常数通常在秒级至分钟级，在小时级的宏观能量调度模型中，忽略暂态过渡过程采用分段恒定功率假设是合理的。该假设在有效降低模型求解维度的同时，确保了日间能量平衡计算的准确性。

%各时段内功率恒定。题目规定"运行分析时段为1小时"，采用分段恒定功率假设，即每个时段内风电、光伏出力和负荷功率保持不变。该假设简化了计算而不影响日能量平衡的准确性。

\textbf{假设五：}不计园区功率损耗。问题一明确给出此条件，线路损耗、变压器损耗等忽略不计。在灵敏度分析中可通过引入3\%--5\%损耗系数评估该假设对结果的影响。

\textbf{假设六：}储能初末荷电状态相同。以日为运行周期，储能系统初始和终末荷电状态均设为额定容量的50\%，确保模型可重复性和不同方案经济性比较的公平性。
